\section{Compute, Bitter Lesson y el ciclo de vida del sesgo inductivo}

Ahora entramos en la línea principal de Hyung Won. En la primera mitad las curvas se trazan en función de la escala del modelo; en la segunda, en función del tiempo histórico. El eje horizontal común sigue siendo el cómputo disponible y la capacidad de los algoritmos. El juicio central de Hyung Won es que la caída sostenida del costo del cómputo, junto con el associated scaling, es la fuerza dominante de la investigación en IA. Por eso, al diseñar una estructura no basta con ver si hoy es más precisa; también hay que ver si limitará seguir ampliando la escala mañana.

\subsection{Por qué el compute per dollar se toma como fuerza dominante}

Esta sección empieza por la evidencia. El eje vertical de la figura no es el número de transistores, sino la capacidad de cómputo que puede comprarse con una cantidad fija de dinero; el eje horizontal abarca más de un siglo. A partir de esta tendencia de largo plazo, Hyung Won propone una estrategia de investigación: no competir contra una tendencia externa fuerte y persistente, sino diseñar métodos que la aprovechen. Al mismo tiempo, hay que incluir el hardware, los algoritmos, la precisión numérica y la eficiencia del software en el «cómputo efectivo»; de lo contrario, se confunde la métrica de un solo componente con la trayectoria completa de la investigación.

\lecturefigure{hyung-slide-007.jpg}{La capacidad de cómputo que se obtiene a costo fijo crece a largo plazo, aproximadamente diez veces cada cinco años}{Hyung Won Chung official deck, p.~7}

Si la capacidad de cómputo crece aproximadamente diez veces cada cinco años, puede escribirse
\[
C(t)=C_0\,10^{t/5},
\]
donde $C(t)$ es la cantidad de cómputo que puede comprarse con un presupuesto fijo $t$ años después del punto de partida, y $C_0$ es la cantidad de cómputo en el punto de partida. Esta expresión solo aproxima la figura de la clase y no garantiza que el mismo exponente se mantenga en el futuro; la energía, la fabricación, la interconexión, el capital y los algoritmos pueden modificar la tasa de crecimiento efectiva.

\begin{knowledgebox}{Lectura de la figura: compute availability no es lo mismo que Moore's law}
En la sesión de preguntas y respuestas, el ponente distingue explícitamente entre la densidad de transistores y el cómputo disponible. Las GPU, la baja precisión, los aceleradores especializados, la pila de software y el codiseño con el hardware pueden seguir mejorando el rendimiento del entrenamiento, aunque el escalamiento tradicional de los transistores se desacelere. La pregunta pertinente es «cuánto cómputo de entrenamiento efectivo puede realizarse con un costo dado», no la métrica de un solo componente.
\end{knowledgebox}

\teachervoice{Entre 00:38:24 y 00:39:27 la formulación es «no compitas con esta tendencia; aprovéchala tanto como puedas». Es una heurística para elegir líneas de investigación, no una garantía de crecimiento exponencial ilimitado; después de 01:15:17, en la sesión de preguntas y respuestas, el ponente menciona por iniciativa propia que la energía y los límites físicos podrían convertirse en nuevos cuellos de botella.}

\subsection{Bitter Lesson: imponerle menos a la máquina la forma de pensar que creemos entender}

La sección anterior mostró que la escala es una tendencia fuerte; esta explica por qué el exceso de estructura diseñada a mano entra en conflicto con ella. Los investigadores suelen codificar en el modelo su propia explicación del pensamiento humano; con poco cómputo, estas estructuras pueden parecer atajos, pero limitan la exploración de otras representaciones por parte del modelo. La versión de la Bitter Lesson que se presenta en clase es esta: el progreso a largo plazo ha surgido una y otra vez de métodos más generales, con supuestos más débiles, sumados a más datos y más cómputo.

\lecturefigure{hyung-slide-008.jpg}{«Enseñar a la máquina a pensar como creemos que pensamos» puede imponer una estructura equivocada}{Hyung Won Chung official deck, p.~8}

\lecturefigure{hyung-slide-009.jpg}{Bitter Lesson: métodos más generales, supuestos de modelado más débiles y mayor escala}{Hyung Won Chung official deck, p.~9}

\begin{importantbox}{Qué es el inductive bias}
El inductive bias (sesgo inductivo) es el supuesto estructural por el cual, con datos limitados, un modelo prefiere cierta clase de funciones, representaciones o formas de interacción. La localidad de la convolución, la separación de parámetros entre encoder y decoder, la attention bidireccional y ciertas invariancias específicas son sesgos. Un sesgo puede mejorar la eficiencia muestral, pero también excluye parte de las soluciones; su valor depende de los datos, el cómputo, la tarea y las restricciones de despliegue del momento.
\end{importantbox}

\teachervoice{Entre 00:39:41 y 00:41:43, Hyung Won dice que los investigadores suelen intentar modelar «cómo creemos que pensamos», sin entender realmente los mecanismos cognitivos de bajo nivel. La postura de la clase es deliberadamente tajante; estas notas la conservan como una advertencia de diseño, no como una demostración de que toda estructura previa sea errónea.}

\subsection{Cuanta menos estructura, más scalable, pero no necesariamente mejor hoy}

Si se compara según la cantidad de estructura, los métodos con más estructura suelen partir de un nivel más alto en el intervalo de poco cómputo, pero pueden saturarse antes; los métodos con menos estructura son difíciles de entrenar al principio, pero dan más libertad al modelo y, al crecer la escala, pueden superar a los primeros. Aquí ``scalable'' significa que la curva de crecimiento puede seguir mejorando, no solo que el método pueda ejecutarse en más GPU.

\lecturefigure{hyung-slide-010.jpg}{Más estructura mejora el desempeño con poco cómputo, pero puede reducir la escalabilidad a largo plazo}{Hyung Won Chung official deck, p.~10}

\begin{knowledgebox}{Lectura de la figura: comparar el punto de cruce en lugar de calificar la estructura como buena o mala}
A la izquierda del presupuesto actual, el método azul puede ser claramente mejor; el método rojo solo lo supera con más compute. La decisión de investigación depende de dónde esté el punto de despliegue, de si puede asumirse el costo de exploración y de si en el futuro realmente se llegará al punto de cruce. Lo que la figura respalda es una regime-dependent choice, no «elegir siempre la menor estructura posible».
\end{knowledgebox}

\lecturefigure{hyung-slide-013.jpg}{En el compute regime actual conviene elegir un sesgo que funcione, sin olvidar retirarlo en el futuro}{Hyung Won Chung official deck, p.~13}

\lecturefigure{hyung-slide-014.jpg}{Dados el cómputo, los datos, los algoritmos y la arquitectura, existe un inductive bias óptimo de manera temporal}{Hyung Won Chung official deck, p.~14}

Si la intensidad de la estructura se denota por $b$ y el estado actual de los recursos por $r$, el desempeño se escribe como
\[
b^{\star}(r)=\arg\max_b\; P(b;r).
\]
Aquí $P(b;r)$ es el desempeño del sesgo $b$ en el estado de recursos $r$, y $b^{\star}(r)$ es la estructura óptima actual. A medida que cambian el compute, los data y los algorithm que contiene $r$, $b^{\star}$ también se desplaza; por eso, una decisión de arquitectura que alguna vez fue correcta no debe congelarse para siempre.

\teachervoice{La salvedad central entre 00:42:28 y 00:43:51 es que no se puede esperar desde el principio al método más general, porque podría no funcionar en absoluto con el presupuesto actual. Lo correcto es «incorporar ahora el sesgo útil y dejar registrado explícitamente cuándo revisarlo y retirarlo en el futuro».}

\subsection{Por qué los incentivos de la investigación favorecen agregar estructura y no eliminarla}

El ciclo de vida de la estructura también tiene causas sociales. Un módulo, un término de pérdida o un mecanismo nuevo se convierte fácilmente en la contribución de un artículo; demostrar que un módulo complejo puede eliminarse a mayor escala suele parecer poco novedoso. Así, la comunidad acumula supuestos diseñados para un regime anterior, pero carece de una práctica sistemática de eliminación. Esta sección considera los incentivos de la investigación como parte de la evolución de las arquitecturas, porque qué experimentos reciben recursos y qué resultados negativos se publican influye directamente en cuánto tiempo se conservan los sesgos antiguos.

\lecturefigure{hyung-slide-015.jpg}{El método que es mejor a largo plazo suele ser peor en el presente}{Hyung Won Chung official deck, p.~15}

\begin{warningbox}{Una ventaja de corto plazo en el leaderboard no equivale a valor de investigación a largo plazo}
Si un método solo lleva la delantera con modelos pequeños, pocos datos o un benchmark específico, hay que seguir preguntando: ¿la ventaja proviene de un prior más fuerte? ¿Se satura al aumentar la escala? ¿La complejidad de implementación dificulta entrenamientos más grandes? Estas preguntas no pueden responderse con la clasificación de un único punto actual.
\end{warningbox}

\teachervoice{Entre 00:43:51 y 00:44:47, Hyung Won afirma sin rodeos que los incentivos de publicación fomentan agregar estructura, pero rara vez recompensan eliminarla. También dice que el método que es mejor a largo plazo casi inevitablemente parece peor hoy, porque un sistema de aprendizaje con más libertad es más caótico en el intervalo de pocos recursos.}

\lecturefigure{hyung-slide-016.jpg}{Resumen de la etapa: la fuerza dominante es el cómputo más barato y el associated scaling}{Hyung Won Chung official deck, p.~16}

Esta diapositiva de resumen no niega todas las demás fuerzas; anuncia el sistema de coordenadas del análisis siguiente: explicar, una por una, las estructuras adicionales de la arquitectura Transformer como inductive bias propios de ciertas condiciones históricas, y preguntar si la escala actual todavía las necesita. Así, el diagrama de arquitectura deja de ser una clasificación estática y se convierte en un objeto sobre el que se puede razonar a lo largo de la scale trajectory.

\subsection{Resumen de la sección}

Esta sección estableció el marco de juicio de Hyung Won: la caída sostenida del costo del cómputo da a los métodos generales un gran potencial a largo plazo; el sesgo inductivo mejora la eficiencia cuando hay pocos recursos, pero puede limitar el escalamiento posterior. La conclusión correcta no es «cuanta más estructura, mejor» ni «cuanta menos estructura, mejor», sino estimar continuamente el sesgo óptimo del regime actual y dejar abierto el camino para futuros experimentos de eliminación.
